\section{Introduction}
Gravitational waves provide a unique observational window into astrophysical and cosmological phenomena. Since the first direct detection in 2015 \cite{abbott2016observation}, the \ac{LVK} collaboration has observed dozens of compact binary coalescences in the 10 Hz to 10 kHz frequency band \cite{}. Accessing gravitational waves at millihertz frequencies requires a space-based detector to overcome terrestrial seismic and Newtonian noise and to provide sufficiently long interferometric baselines. The \ac{LISA}, a joint ESA–NASA mission scheduled for launch in the mid-2030s, is designed to operate in the 0.1 mHz to 1 Hz band \cite{colpi_lisa_2024}.

The \ac{LISA} constellation consists of three spacecraft arranged in an approximately equilateral triangle with arm lengths of about 2.5 million km. Laser beams exchanged between the spacecraft form a constellation-scale interferometer. Unlike ground-based detectors with effectively equal arms, the LISA arm lengths vary by up to about 1.2\% due to orbital motion, preventing direct cancellation of laser frequency noise. Although cavity pre-stabilisation reduces laser frequency noise substantially, it remains several orders of magnitude above the displacement noise budget required for gravitational-wave detection \cite{}. To suppress this noise, \ac{TDI} has been developed as a post-processing technique that synthesizes virtual equal-arm interferometers by appropriately time-shifting and combining phasemeter measurements.

In addition to laser frequency noise, LISA measurements are affected by clock noise originating from the independent \acp{USO} carried by each spacecraft. Since all phase measurements are referenced to these clocks, timing jitter couples directly into the measured heterodyne beatnotes. Space-qualified \acp{USO} achieve Allan deviations of order $10^{-13}$ at one-second averaging times \cite{weaver2010performance}, which is insufficient for the picometer-level test-mass position sensitivity required by LISA. To mitigate this effect, LISA employs additional clock-noise transfer measurements—implemented via modulation sidebands exchanged between spacecraft.
 
%Experimental tests of \ac{TDI}, including representative inter-spacecraft delays implemented in hardware, have been performed in several laboratory setups. These experiments have demonstrated laser-frequency noise suppression, clock-noise suppression, and electronic delay implementations on the order of seconds. The present work builds upon these efforts by extending such testbeds toward a more integrated hardware-in-the-loop configuration and by enabling the injection of physically consistent gravitational-wave strain signals into the measurement chain.

\ac{TDI} has been extensively studied theoretically [] and validated through numerical simulations []. Such simulations, however, generate the measurements from the same model that the analysis pipeline assumes, and therefore test the implementation of the model rather than the model itself. Testing the model requires data produced independently of it, by a physical realization of the measurement chain. Experimental verification of this kind is still developing.

A major milestone was the \ac{UFLIS}, which demonstrated that second-generation \ac{TDI} combinations can suppress laser frequency noise by ten orders of magnitude [4]. The work presented here is a continuation of the \ac{UFLIS} approach with the addition of independent clock sources for each spacecraft, largely enabled by developments in FPGA technology. Other experiments have also contributed important demonstrations relevant to \ac{TDI} and clock-noise removal. The setup by de Vine et al. demonstrated laser frequency noise suppression at a level of roughly nine orders of magnitude and clock-phase noise suppression of about four orders, although without accounting for realistic photon light-travel times. The LISA on Table experiment in Paris more recently demonstrated laser-noise removal using direct digital synthesis (DDS)-generated signals and electronically imposed delays. The Hexagon experiment in Hannover has provided strong results for clock-noise removal and synchronization.

Here, we present an overview of an ongoing effort to develop a \ac{TDI} testbed that combines realistic laser frequency noise, programmable delay lines, and independent clock sources within a single experimental framework. The testbed focuses on reproducing the dominant noise couplings of the LISA metrology system—specifically, laser frequency noise and clock noise—while preserving the structure imposed by unequal and time-dependent delays. Crucially, the same delay lines that impose the inter-spacecraft light-travel times also allow physically consistent gravitational-wave strain signals to be injected directly into the measurement chain. The testbed can therefore be used not only to demonstrate noise suppression, but to verify that a known gravitational-wave signal survives \ac{TDI} processing of hardware-generated data with the expected amplitude and phase.

An initial purely electrical implementation centered around \ac{FPGA}-based digital delay lines has already been demonstrated in Ferguson et al. [REF]. The developments discussed here extend that architecture by incorporating an optical front-end with physically representative laser noise and independently running spacecraft clocks. This preserves the essential timing, noise, and signal-propagation structure of the LISA one-way measurements, allowing \ac{TDI} combinations and clock-noise calibration to be formed entirely from hardware-generated data. The goal of this work is to realize a laboratory-scale analog of the LISA inter-satellite metrology system—a “miniLISA”—that bridges the gap between idealized software simulations and the flight instrument, and provides a controlled environment for testing the robustness of \ac{TDI} to instrumental non-idealities and realistic gravitational-wave signal injection.

The goal is for both noise sources, once processed through \ac{TDI} and clock-noise calibration, to be reduced below the level required for gravitational-wave signal recovery across the core LISA measurement band.

The article is divided as follows:% in section \ref{}, we give an overview of the \ac{LISA} measurement system, and in section \ref{}, we describe our experimental setup to replicate said measurement system in the lab with a focus on aspects relevant to \ac{TDI} and clock noise removal. In section \ref{Comparison}, analyze the expected performance of the setup. In section \ref{}, we discuss the limitations and future prospects in replicating the \ac{LISA} measurement system in this manner, and in section \ref{}, we conclude.